\reading{IM-5}{Portfolio Risk: Analytical Methods}{STRIKE FIRST}{6}

\srcnote{Philippe Jorion, \textit{Value-at-Risk: The New Benchmark for Managing
Financial Risk}, 3rd Edition (New York, NY: McGraw-Hill, 2007), Chapter 7.
Printed as Chapter 5 in the GARP source volume.}

\begin{imkeybox}[title={The five objectives in this reading}]
\begin{enumerate}[label=\textbf{\alph*.}]
\item Define, calculate, and compare the following portfolio VaR measures: diversified and undiversified portfolio VaR, individual VaR, incremental VaR, marginal VaR, and component VaR.
\item Explain the impact of correlation on portfolio risk.
\item Apply the concept of marginal VaR in making portfolio management decisions.
\item Explain and calculate the risk-minimizing position and the position that maximizes the ratio of expected return to risk.
\item Explain the difference between risk management and portfolio management and describe how to use marginal VaR in portfolio management.
\end{enumerate}
\end{imkeybox}

\begin{imtrapbox}[breakable=false,title={Read the notation before reading the formulas}]
This reading conflates three distinct symbols, and it does so \emph{within itself}.
The crash layer writes incremental VaR as $\text{MVaR}_A \times W_A$ while the
lecture layer writes the same quantity as $\text{MVaR}_A \times P_A$; one slide
uses $w_i$ for a portfolio \emph{weight} and another uses $w_T$ for a
\emph{dollar position} of CAD 15 million. Throughout this volume:
\term{$w_i$} is a weight, \term{$W$} is total portfolio value, and
\term{$x_i$ or $P_i$} is a dollar position. Both of the notes' usages are carried
across where they appear, with the meaning named each time. They are not silently
merged.
\end{imtrapbox}

\basics{Basics: the delta-normal setting}

\begin{itemize}
\item \term{Delta-normal model:} all individual security returns are assumed
normally distributed.
\item \term{Traditional portfolio analysis:} based on variances and covariances,
which is why it is sometimes called the \term{covariance matrix approach}.
\end{itemize}

\figwrap{%
\begin{tikzpicture}[
  hd/.style={draw=FRMNavy,fill=FRMNavy,rounded corners=1pt,inner sep=4pt,
             font=\scriptsize\sffamily\bfseries,text=white,align=center,
             text width=32mm},
  bx/.style={draw=FRMRule,fill=white,rounded corners=1pt,inner sep=3.5pt,
             font=\scriptsize\sffamily,align=center,text width=34mm},
  ar/.style={-{Stealth[length=4pt]},FRMGold,line width=0.9pt}]
\node[hd] (a) at (0,0) {Portfolio VaR measures};
\node[bx,below=5mm of a] (a1) {individual VaR};
\node[bx,below=2.5mm of a1] (a2) {diversified VaR};
\node[bx,below=2.5mm of a2] (a3) {undiversified VaR};
\draw[ar] (a) -- (a1);
\node[hd] (b) at (4.9,0) {VaR tools};
\node[bx,below=5mm of b] (b1) {marginal VaR};
\node[bx,below=2.5mm of b1] (b2) {incremental VaR};
\node[bx,below=2.5mm of b2] (b3) {component VaR};
\draw[ar] (b) -- (b1);
\node[hd] (c) at (9.8,0) {Risk management};
\node[bx,below=5mm of c] (c1) {global minimum risk:\\ all MVaR equal};
\node[bx,below=2.5mm of c1] (c2) {optimal portfolio:\\ all $E_i/\text{MVaR}_i$ equal};
\draw[ar] (c) -- (c1);
\draw[ar] (a.east) -- (b.west);
\draw[ar] (b.east) -- (c.west);
\end{tikzpicture}}%
{The reading in three moves, redrawn from the source notes' own sketch. The first
column \emph{measures} risk, the second \emph{decomposes} it so you can see which
position is responsible, and the third \emph{acts} on it. Every objective in this
reading sits in one of the three columns, and the arrows between them are the
argument of the chapter: you cannot manage what you have not decomposed.}

% ------------------------------------------------------------------
\lo{IM-5 a}{Define, calculate, and compare the following portfolio VaR measures: diversified and undiversified portfolio VaR, individual VaR, incremental VaR, marginal VaR, and component VaR.}

\subsection{Individual VaR}

This is the VaR of an individual position \term{in isolation}.

\begin{imfmlbox}
\[ \text{VaR}_i = Z_c \times \sigma_i \times |P_i| = Z_c \times \sigma_i \times |w_i| \times P \]
\vk{$P$ is portfolio value, $P_i$ the nominal amount invested in position $i$,
$w_i$ its weight, $\sigma_i$ its volatility, and $Z_c$ the confidence multiplier.}
\end{imfmlbox}

\subsection{Diversified and undiversified VaR}

\term{Diversified VaR} is the portfolio VaR, taking into account diversification
benefits between components:
\begin{imfmlbox}
\[ \text{VaR}_p = Z_c \times \sigma_p \times P
   \qquad\text{with}\qquad
   \sigma_p^2 = w_1^2\sigma_1^2 + w_2^2\sigma_2^2 + 2w_1w_2\rho_{12}\sigma_1\sigma_2 \]
\end{imfmlbox}

\term{Undiversified VaR} is simply the sum of the individual VaRs. If correlation
is below 1, diversified VaR is less than undiversified VaR.

\begin{imfmlbox}
\[ \rho = 1:\quad \text{VaR}_p = \text{VaR}_1 + \text{VaR}_2
   \qquad\qquad
   \rho = 0:\quad \text{VaR}_p = \sqrt{\text{VaR}_1^2 + \text{VaR}_2^2} \]
\end{imfmlbox}

\subsection{The equally weighted $N$-asset case}

There are three assumptions behind this formula:

\begin{itemize}
\item Each asset of the portfolio is equally weighted.
\item Each individual position has the same standard deviation of return.
\item Each pair of returns has only one correlation.
\end{itemize}

\begin{imfmlbox}
\[ \sigma_P = \sigma\sqrt{\frac{1}{N} + \left(1 - \frac{1}{N}\right)\rho} \]
\vk{$N$ is the number of positions, $\sigma$ the standard deviation for each
position, and $\rho$ the correlation for each pair of positions.}
\end{imfmlbox}

\begin{imexbox}[title={Five positions, one correlation}]
A portfolio has five positions of \$2 million each. The standard deviation of the
returns is 30\% for each position. The correlation between each pair of returns
is 0.2. Calculate the VaR using a $z$-score of 2.33.
\begin{align*}
\sigma_P &= 30\%\sqrt{\tfrac{1}{5} + \left(1 - \tfrac{1}{5}\right)0.2}
  = 30\%\sqrt{0.2 + 0.16} = 30\%\sqrt{0.36} = 18\% \\[2pt]
\text{VaR}_p &= Z_c \times \sigma_p \times V = (2.33)(18\%)(\$10\text{ million})
  = \mathbf{\$4{,}194{,}000}
\end{align*}
\end{imexbox}

\begin{imtrapbox}[breakable=false,title={The \$10 million is not given}]
Five positions of \$2 million each. The stem gives the \emph{position} size and
the answer needs the \emph{portfolio} size. Substituting \$2 million gives
\$838,800, which is a real and plausible-looking number. This is the standard
intermediate-result trap in this family of questions.
\end{imtrapbox}

\subsection{Marginal VaR}

\begin{imdefbox}[title={Marginal VaR}]
The change in portfolio VaR resulting from taking an \term{additional dollar} of
exposure to a given component. It is also the partial derivative with respect to
the component position.
\end{imdefbox}

\begin{imfmlbox}
\[ \text{MVaR}_i = \frac{\partial\,\text{VaR}_p}{\partial(\text{monetary investment in } i)}
  = Z_c\frac{\partial \sigma_p}{\partial w_i}
  = Z_c\frac{\operatorname{cov}(R_i,R_p)}{\sigma_p} \]
\vspace{-10pt}
\[ \text{MVaR}_i = \frac{\text{VaR}_p}{\text{portfolio value}} \times \beta_i \]
\vk{$\beta_i = \operatorname{cov}(R_i,R_p)/\sigma_p^2$. The second form is the one
to reach for when the question gives a beta, and the first when it gives a
covariance matrix.}
\end{imfmlbox}

\begin{imexbox}[title={Marginal VaR from a beta}]
Portfolio X has a VaR of EUR 400,000. It holds four equally weighted assets, each
valued at EUR 1,000,000. Asset A has a beta of 1.2.
\[ \text{MVaR}_A = \frac{\text{VaR}_P}{\text{portfolio value}} \times \beta_A
  = \frac{400{,}000}{4{,}000{,}000} \times 1.2 = \mathbf{0.12} \]
Portfolio VaR will change by 0.12 for each euro change in Asset A.
\end{imexbox}

\subsection{Incremental VaR}

\begin{imdefbox}[title={Incremental VaR}]
The change in VaR owing to a \term{new position}. It differs from marginal VaR in
that the amount added or subtracted can be large.
\end{imdefbox}

Incremental VaR requires a \term{full revaluation} of the portfolio VaR with the
new trade:

\begin{imfmlbox}
\[ \text{Incremental VaR}_A = \text{VaR}_{P+A} - \text{VaR}_P
   \qquad\qquad
   \text{Incremental VaR}_A \approx \text{MVaR}_A \times P_A \]
\vk{$P_A$ is the \emph{dollar} size of position $A$. The crash layer of the source
notes writes this same quantity as $W_A$; it is a position, not the portfolio's
total value.}
\end{imfmlbox}

The main drawback of the full-revaluation approach is that it requires a full
revaluation of the portfolio VaR with the new trade, which is time-consuming for
large portfolios. \term{If VaR is decreased, the new trade is risk-reducing or is
a hedge}; otherwise the new trade is risk-increasing.

\begin{imexbox}[title={Incremental VaR on an uncorrelated pair}]
Assets A and B have volatilities of 6\% and 14\%, with \$4 million and \$2 million
invested. They are uncorrelated. Compute the incremental VaR for an increase of
\$10,000 in Asset A, at $z = 1.65$.

Per-dollar covariances of each risk factor with the portfolio (in \$m):
\[
\begin{bmatrix}\operatorname{cov}(R_A,R_P)\\ \operatorname{cov}(R_B,R_P)\end{bmatrix}
=\begin{bmatrix}0.06^2 & 0\\ 0 & 0.14^2\end{bmatrix}
 \begin{bmatrix}\$4\\ \$2\end{bmatrix}
=\begin{bmatrix}0.0144\\ 0.0392\end{bmatrix}
\]
The portfolio variance in dollar terms is
$4(0.0144) + 2(0.0392) = 0.136$, so $\sigma_P = \sqrt{0.136}$.
\begin{align*}
\text{MVaR}_A &= 1.65 \times \frac{0.0144}{\sqrt{0.136}} = 0.064428 \\
\text{MVaR}_B &= 1.65 \times \frac{0.0392}{\sqrt{0.136}} = 0.175388
\end{align*}
Since the two assets are uncorrelated, the incremental VaR of an additional
\$10,000 in Position A is $\$10{,}000 \times 0.064428 = \mathbf{\$644.28}$.
\end{imexbox}

\subsection{Component VaR}

\begin{imdefbox}[title={Component VaR}]
A partition of the portfolio VaR that indicates \term{how much the portfolio VaR
would change approximately if the given component were deleted}. By construction,
component VaRs sum to portfolio VaR.
\end{imdefbox}

\begin{imfmlbox}
\[ \text{CVaR}_i = (\text{MVaR}_i)(w_i \times P) = \text{VaR}_p \times \beta_i \times w_i
   = (Z_c \sigma_i w_i P) \times \rho_i = \text{VaR}_i \times \rho_i \]
\vk{the last form is the one worth memorising: \term{component VaR is individual
VaR times the correlation with the portfolio}. Here $w_i$ is a weight and $P$ the
portfolio value.}
\end{imfmlbox}

\begin{imexbox}[title={Component VaR from a correlation}]
A portfolio is valued at CAD 248 million and contains CAD 15 million in stock T.
The annualised standard deviations of the portfolio and of stock T are 16\% and
13\%. The correlation between them is 0.45. Using a 1-year 95\% VaR with normal
returns, what is the component VaR of stock T?

\medskip
\begin{tabular}{@{}ll@{}}
A & CAD 0.096 million \\
B & CAD 1.444 million \\
C & CAD 2.041 million \\
D & CAD 3.948 million \\
\end{tabular}

\medskip
$\text{VaR}_T = 15 \times 0.13 \times 1.645 = \text{CAD } 3.208$ million, so
\[ \text{CVaR}_T = \rho_{T,p} \times \text{VaR}_T = 0.45 \times 3.208
   = \mathbf{\text{CAD } 1.444\text{ million}} \quad\text{(B)} \]
\end{imexbox}

\begin{imtrapbox}[breakable=false,title={Every distractor here is a real quantity}]
CAD 3.948 is the individual VaR computed with the \emph{portfolio's} 16\%
volatility instead of the stock's 13\%. CAD 3.208 --- the individual VaR itself
--- is the classic intermediate-result answer and would be offered on a different
day. Note also that this stem uses $w_T$ for the \emph{dollar value} of stock T,
not a weight. The portfolio's CAD 248 million is never used: it is a red herring.
\end{imtrapbox}

\begin{imexbox}[title={Component and marginal VaR from a table}]
\begin{center}\small
\begin{tabular}{@{}lrrr@{}}
\toprule
Asset & Position value (\$000) & Return s.d.\ (\%) & Beta \\
\midrule
A & 400 & 3.60 & 0.5 \\
B & 600 & 8.63 & 1.2 \\
Portfolio & 1{,}000 & 5.92 & 1 \\
\bottomrule
\end{tabular}
\end{center}
\medskip
Calculate the component VaR of asset A and the marginal VaR of asset B at 95\%.

\medskip
\begin{tabular}{@{}ll@{}}
A. & USD 21,773 and 0.1306 \\
B. & USD 21,773 and 0.1169 \\
C. & USD 19,477 and 0.1169 \\
D. & USD 19,477 and 0.1306 \\
\end{tabular}

\medskip
\begin{align*}
\text{VaR}_p &= \alpha \times \sigma_p \times P = 1.645 \times 0.0592 \times 1{,}000{,}000 = \text{USD } 97{,}384\\
\text{CVaR}_A &= 97{,}384 \times 0.5 \times \tfrac{400}{1000} = \text{USD } \mathbf{19{,}477}\\
\text{MVaR}_B &= 97{,}384 \times 1.2 / 1{,}000{,}000 = \mathbf{0.1169}
\end{align*}
Answer: \textbf{C}.
\end{imexbox}

\begin{imtrapbox}[breakable=false,title={A two-by-two grid}]
The four options are every combination of two right answers and two wrong ones.
Getting the component VaR right and the marginal VaR wrong lands on a real
offered answer, and so does the reverse. Neither half can be checked against the
other --- which is the point of building options this way.
\end{imtrapbox}

\subsection{Computing portfolio VaR with the covariance matrix}

\begin{imexbox}[title={A two-asset 10-day VaR}]
A two-asset portfolio holds \$8 million in asset A and \$4 million in asset B. The
correlation is 0.6; daily standard deviations are 1.5\% and 2.0\%. What is the
10-day VaR at 99\% ($\alpha = 2.33$)?

\medskip
The covariance matrix, with $\operatorname{cov}_{12} = 0.6(0.015)(0.02) = 0.00018$:
\[
\begin{bmatrix}0.000225 & 0.00018\\ 0.00018 & 0.0004\end{bmatrix}
\]
Post-multiplying by the position vector and then pre-multiplying:
\[
\begin{bmatrix}8 & 4\end{bmatrix}
\begin{bmatrix}0.000225 & 0.00018\\ 0.00018 & 0.0004\end{bmatrix}
= \begin{bmatrix}0.00252 & 0.00304\end{bmatrix},
\qquad
\begin{bmatrix}0.00252 & 0.00304\end{bmatrix}\begin{bmatrix}8\\4\end{bmatrix} = 0.03232
\]
so the \emph{dollar} standard deviation is $\sqrt{0.03232} = 0.1798$ (\$m), and
\[ \text{VaR}_P = \sigma_P \times \alpha \times \sqrt{t}
   = 0.1798 \times 2.33 \times \sqrt{10} = \mathbf{\$1.3248\text{ million}} \]
\end{imexbox}

\begin{imtrapbox}[breakable=false,title={A dollar sigma, not a percentage sigma}]
The source notes write this last line as $\text{VaR}_P = \sigma_P\alpha\sqrt{x}$.
Two things are hidden in that. The $x$ is the \emph{holding period} in days, not a
variable quantity; and $\sigma_P = 0.1798$ is measured in \emph{millions of
dollars}, because the position vector was in millions. Multiplying it by the
portfolio value again --- the reflex from the previous example --- overstates the
answer by a factor of twelve.
\end{imtrapbox}

% ------------------------------------------------------------------
\lo{IM-5 b}{Explain the impact of correlation on portfolio risk.}

Correlation sets where portfolio risk sits between two hard bounds. At $\rho=1$
the individual VaRs simply add; at $\rho=0$ they add in quadrature; and as $N$
grows the equally weighted portfolio's volatility does not fall to zero but to a
floor of $\sigma\sqrt{\rho}$.

\figwrap{%
\begin{tikzpicture}
\begin{axis}[frmaxis,
  width=11.4cm, height=5.6cm,
  xmin=1, xmax=50, ymin=0, ymax=0.32,
  xlabel={number of equally weighted positions $N$},
  ylabel={portfolio volatility $\sigma_P$},
  ytick={0,0.05,0.10,0.134,0.1897,0.30},
  yticklabels={0,5\%,10\%,13.4\%,19.0\%,30\%},
  legend style={font=\scriptsize\sffamily, draw=FRMRule, fill=white,
    at={(0.5,-0.26)}, anchor=north, legend columns=-1, column sep=10pt},
]
\addplot[FRMRed, very thick, domain=1:50, samples=120]
  {0.30*sqrt(1/x + (1-1/x)*0.4)};
\addlegendentry{$\rho = 0.4$}
\addplot[FRMNavy, very thick, domain=1:50, samples=120]
  {0.30*sqrt(1/x + (1-1/x)*0.2)};
\addlegendentry{$\rho = 0.2$}
\addplot[FRMGreen, very thick, domain=1:50, samples=120]
  {0.30*sqrt(1/x)};
\addlegendentry{$\rho = 0$}
\draw[FRMRed, dashed, thin] (axis cs:1,0.1897) -- (axis cs:50,0.1897);
\draw[FRMNavy, dashed, thin] (axis cs:1,0.1342) -- (axis cs:50,0.1342);
\node[font=\scriptsize\sffamily, fill=white, inner sep=1.5pt, anchor=west,
      text=FRMRed] at (axis cs:22,0.215) {floor $=\sigma\sqrt{\rho}$};
\end{axis}
\end{tikzpicture}}%
{Why correlation, not the number of holdings, decides how much diversification
you get. Each curve is
$\sigma_P = \sigma\sqrt{1/N + (1-1/N)\rho}$ with $\sigma = 30\%$. Only the
$\rho = 0$ case decays towards zero. With $\rho = 0.2$ the volatility is already
within a whisker of its $13.4\%$ floor by twenty holdings, and adding a
twenty-first position buys almost nothing. The dashed lines are $\sigma\sqrt{\rho}$
--- the limit as $N \to \infty$, which is where $1/N$ vanishes and only the
correlation term survives.}

\begin{imkeybox}
This is the same arithmetic as the five-position example above: $N = 5$,
$\rho = 0.2$ and $\sigma = 30\%$ gives $18\%$, which is already two-thirds of the
way from $30\%$ down to the $13.4\%$ floor.
\end{imkeybox}

% ------------------------------------------------------------------
\lo{IM-5 c}{Apply the concept of marginal VaR in making portfolio management decisions.}

Over time, risk managers discovered that they could use the VaR process for
active risk management. A typical question is: \term{which position should I
alter to modify my VaR most effectively?}

The three tools answer three different versions of that question:

\renewcommand{\arraystretch}{1.25}
\noindent\begin{tabularx}{\textwidth}{@{}>{\raggedright\arraybackslash}p{2.7cm}XX@{}}
\toprule
\textbf{\sffamily\small Measure} & \textbf{\sffamily\small Question it answers} & \textbf{\sffamily\small Size of the change} \\
\midrule
Marginal VaR & What happens if I add one more dollar here? & Infinitesimal --- a derivative, a slope. \\
Incremental VaR & What happens if I put this whole new trade on, or take this whole position off? & Large. Requires full revaluation. \\
Component VaR & Of the risk I already have, how much is this position responsible for? & The current position. Sums to portfolio VaR. \\
\bottomrule
\end{tabularx}
\renewcommand{\arraystretch}{1}
\figcap{The three tools are not three ways of saying the same thing. They differ
by the \emph{size of the change} being contemplated, and component VaR is the
only one of the three that partitions the existing portfolio.}

\gapbadge
\srcnote{The notes carry GARP's Figure 5.2 --- the incremental-VaR revaluation
flow --- but not Figure 5.4, which puts all four measures on a single curve. It is
reconstructed below from GARP's own reported quantities for the two-currency
portfolio, and every value on it has been recomputed.}

\begin{imgapbox}[title={What the objective asks for}]
The four measures are not four formulas. They are four different readings of
\emph{one} curve: the portfolio VaR as a function of how much you hold.
\end{imgapbox}

\figwrap{%
\begin{tikzpicture}
\begin{axis}[frmaxis,
  width=11.6cm, height=6.4cm,
  xmin=-0.35, xmax=2.0, ymin=0, ymax=470,
  xlabel={position in the euro (\$ million)},
  ylabel={portfolio VaR (\$000)},
  xtick={0,0.5,1,1.5,2}, ytick={0,105.6,165,257.7,350},
  yticklabels={0,105.6,165.0,257.7,350},
  legend style={font=\scriptsize\sffamily, draw=FRMRule, fill=white,
    at={(0.5,-0.24)}, anchor=north, legend columns=-1, column sep=10pt},
]
\addplot[FRMNavy, very thick, domain=-0.35:2.0, samples=140]
  {sqrt(27225 + 3.9*x + 39200*x^2)};
\addlegendentry{portfolio VaR}
\addplot[FRMRed, thick, dashed, domain=-0.1:1.6, samples=2]
  {257.74 + 152.1*(x-1)};
\addlegendentry{tangent at the current \$1m position}
\addplot[only marks, mark=*, mark size=2.2pt, FRMNavy] coordinates {(1,257.74)};
% incremental VaR bracket at x=0
\draw[FRMGreen, line width=1pt, {Stealth[length=4pt]}-{Stealth[length=4pt]}]
  (axis cs:0,165) -- (axis cs:0,257.74);
\draw[black!35, thin, dotted] (axis cs:0,257.74) -- (axis cs:1,257.74);
\draw[FRMGold, line width=1pt, {Stealth[length=4pt]}-{Stealth[length=4pt]}]
  (axis cs:0.13,105.6) -- (axis cs:0.13,257.74);
\node[font=\scriptsize\sffamily, fill=white, inner sep=1.5pt, anchor=west,
      text=FRMGreen] at (axis cs:0.21,200) {incremental VaR \$92.7k};
\node[font=\scriptsize\sffamily, fill=white, inner sep=1.5pt, anchor=west,
      text=FRMGold] at (axis cs:0.21,150) {component VaR \$152.1k};
\node[font=\scriptsize\sffamily, fill=white, inner sep=1.5pt, anchor=east,
      text=FRMRed] at (axis cs:1.95,330) {slope $=$ marginal VaR $=0.1521$};
\draw[-{Stealth[length=4pt]}, FRMRed, thin] (axis cs:1.28,322) -- (axis cs:1.12,285);
\node[font=\scriptsize\sffamily, fill=white, inner sep=1.5pt, anchor=south,
      text=FRMNavy] at (axis cs:-0.02,35) {best hedge};
\draw[-{Stealth[length=4pt]}, FRMNavy, thin] (axis cs:-0.02,60) -- (axis cs:0,158);
\end{axis}
\end{tikzpicture}}%
{All four measures on one curve, reconstructed for GARP's two-currency portfolio
(\$2m CAD, \$1m EUR, portfolio VaR \$257,738). \term{Marginal VaR} is the slope of
the tangent at the current position, $0.1521$ per dollar. \term{Component VaR},
\$152.1k, is measured down the \emph{tangent} to a zero euro position.
\term{Incremental VaR}, \$92.7k, is measured down the \emph{curve} to the same
place. The gap between them --- nearly \$60k --- is the whole lesson: component
VaR is a linear approximation to incremental VaR, and it is a poor one here
because the euro is a third of a two-asset portfolio. The \term{best hedge} is
where the curve bottoms out, which for this portfolio is a net zero euro
position.}

\begin{imtrapbox}[breakable=false,title={Down the tangent, or down the curve}]
Component VaR and incremental VaR answer the same question --- what happens if
this position goes away --- and give different numbers, \$152.1k against \$92.7k.
Component VaR travels along the straight line; incremental VaR along the real
curve. The approximation improves when the component is \emph{small} relative to
the portfolio and degrades when it is large. A question that calls them equal, or
that offers the component figure when it asked for the incremental one, is using
exactly this gap.
\end{imtrapbox}

% ------------------------------------------------------------------
\lo{IM-5 d}{Explain and calculate the risk-minimizing position and the position that maximizes the ratio of expected return to risk.}

\subsection{From risk measurement to risk management}

Taking \emph{only risk} into consideration leads to the \term{global minimum}.
A manager can lower a portfolio's VaR by:

\begin{itemize}
\item lowering the position with the \term{highest} marginal VaR;
\item adding to the position with the \term{lowest} marginal VaR.
\end{itemize}

\begin{imkeybox}[title={At the global minimum}]
All marginal VaRs must be equal, and therefore all betas must be equal:
\[ \text{MVaR}_i = \frac{\text{VaR}}{W} \times \beta_i = \text{constant} \]
\end{imkeybox}

\gapbadge
\srcnote{The spine's verb is \emph{explain and calculate}. The notes explain both
conditions and calculate neither. Both worked positions below are GARP's own
(Chapter 5, Tables 5.4 and 5.5) for the two-currency portfolio, and every figure
in them has been independently recomputed.}

\begin{imexbox}[title={Calculating the risk-minimizing position}]
\begin{center}\small
\begin{tabular}{@{}lrrcrr@{}}
\toprule
& \multicolumn{2}{c}{\textbf{Original}} & & \multicolumn{2}{c}{\textbf{Risk-minimizing}} \\
\cmidrule(r){2-3}\cmidrule(l){5-6}
Asset & Position $w_i$ & MVaR & & Position $w_i$ & MVaR \\
\midrule
CAD & 66.67\% & 0.0528 & & 85.21\% & 0.0762 \\
EUR & 33.33\% & 0.1521 & & 14.79\% & 0.0762 \\
\midrule
Diversified VaR & \multicolumn{2}{c}{\$257,738} & & \multicolumn{2}{c}{\$228,462} \\
Volatility & \multicolumn{2}{c}{15.62\%} & & \multicolumn{2}{c}{13.85\%} \\
\bottomrule
\end{tabular}
\end{center}
\medskip
The marginal VaR is higher for the EUR, so \term{the EUR position is cut first
while adding to the CAD}. At the final position the marginal VaRs are equal and,
equivalently, the betas are both 1.000:
\[ \text{MVaR}_i = \frac{\$228{,}462}{\$3{,}000{,}000} \times 1.000 = 0.0762 \]
The EUR weight falls from 33.33\% to 14.79\% and portfolio volatility falls from
15.62\% to 13.85\%.
\end{imexbox}

\subsection{From risk management to portfolio management}

Taking \emph{return and risk} into consideration leads to the \term{optimal
portfolio}: maximise the Sharpe ratio, written here with VaR in the denominator.

\begin{imfmlbox}
\[ \frac{\overline{R_p} - \overline{R_f}}{\text{VaR}_p}
   \qquad\text{maximised when}\qquad
   \frac{R_i - R_f}{\text{MVaR}_i} = \frac{R_j - R_f}{\text{MVaR}_j}
   = \frac{R_i - R_f}{\beta_i} = \text{constant} \]
\end{imfmlbox}

\begin{itemize}
\item \term{Lower} the allocation to the position with the \term{lowest} excess
expected return to marginal VaR.
\item \term{Add} to the position with the \term{highest} excess expected return to
marginal VaR.
\end{itemize}

\begin{imtrapbox}[breakable=false,title={A subscript the source notes lost}]
The summary slide in the notes prints the optimality condition with
$\text{MVaR}_i$ in \emph{both} denominators, so it reads as an identity that is
trivially true. It should be $\text{MVaR}_i$ on the left and $\text{MVaR}_j$ on
the right --- the whole content of the condition is that the ratio is the same
\emph{across different assets}.
\end{imtrapbox}

\begin{imexbox}[title={Calculating the return-optimizing position}]
Assume expected returns of 8\% for CAD and 5\% for EUR, in excess of the risk-free
rate.
\begin{center}\small
\begin{tabular}{@{}lrrrcrrr@{}}
\toprule
& \multicolumn{3}{c}{\textbf{Original}} & & \multicolumn{3}{c}{\textbf{Optimal}} \\
\cmidrule(r){2-4}\cmidrule(l){6-8}
Asset & $E_i$ & $w_i$ & $E_i/\beta_i$ & & $w_i$ & $\beta_i$ & $E_i/\beta_i$ \\
\midrule
CAD & 8.00\% & 66.67\% & 0.1301 & & 90.21\% & 1.038 & 0.0771 \\
EUR & 5.00\% & 33.33\% & 0.0282 & & \phantom{0}9.79\% & 0.649 & 0.0771 \\
\midrule
Diversified VaR & \multicolumn{3}{c}{\$257,738} & & \multicolumn{3}{c}{\$230,720} \\
Volatility & \multicolumn{3}{c}{15.62\%} & & \multicolumn{3}{c}{13.98\%} \\
Expected return & \multicolumn{3}{c}{7.00\%} & & \multicolumn{3}{c}{7.71\%} \\
Sharpe ratio & \multicolumn{3}{c}{0.448} & & \multicolumn{3}{c}{0.551} \\
\bottomrule
\end{tabular}
\end{center}
\medskip
The ratio $E_i/\beta_i$ is $0.08/0.615 = 0.1301$ for the CAD, greater than
$0.05/1.770 = 0.0282$ for the EUR, so \term{the CAD position should be increased}.
At the optimum the two ratios are identical at
$0.08/1.038 = 0.05/0.649 = 0.0771$, which is the signal that there is no reason to
deviate further. The portfolio expected return is
$0.9021(8\%) + 0.0979(5\%) = 7.71\%$ and the Sharpe ratio rises from 0.448 to
\textbf{0.551}.
\end{imexbox}

\begin{imkeybox}[title={A check that always works}]
At both the minimum-risk and the optimal portfolio, the \emph{weighted average
beta must equal one}, because every asset's beta is measured against the portfolio
itself. Original: $0.6667(0.615) + 0.3333(1.770) = 1.000$. Optimal:
$0.9021(1.038) + 0.0979(0.649) = 1.000$. If a candidate answer's weights and betas
do not reconcile to one, the answer is wrong before any other check.
\end{imkeybox}

\figwrap{%
\begin{tikzpicture}
\begin{axis}[frmaxis,
  width=11.4cm, height=6.0cm,
  xmin=12.8, xmax=17.0, ymin=6.2, ymax=8.6,
  xlabel={risk --- portfolio volatility (\%)},
  ylabel={expected excess return (\%)},
  xtick={13.85,15.62}, xticklabels={13.85,15.62},
  ytick={7.00,7.56,7.71}, yticklabels={7.00,7.56,7.71},
]
\addplot[FRMNavy, very thick, smooth, domain=13.85:17, samples=60]
  {7.556 + 1.02*sqrt(x-13.85)};
\addplot[FRMRule, very thick, smooth, domain=13.85:16.4, samples=60]
  {7.556 - 1.02*sqrt(x-13.85)};
\addplot[FRMGold, thick, dashed, domain=0:17, samples=2] {0.5512*x};
\addplot[only marks, mark=*, mark size=2.2pt, FRMRed]
  coordinates {(15.62,7.00)};
\addplot[only marks, mark=square*, mark size=2pt, FRMNavy]
  coordinates {(13.85,7.556)};
\addplot[only marks, mark=*, mark size=2.4pt, FRMGreen]
  coordinates {(13.98,7.71)};
\node[font=\scriptsize\sffamily, fill=white, inner sep=1.5pt, anchor=east,
      text=FRMRed] at (axis cs:15.46,6.82) {initial portfolio, SR $0.448$};
\draw[-{Stealth[length=4pt]}, FRMRed, thin] (axis cs:15.50,6.86) -- (axis cs:15.60,6.97);
\node[font=\scriptsize\sffamily, fill=white, inner sep=1.5pt, anchor=east,
      text=FRMNavy] at (axis cs:13.80,7.40) {minimum risk};
\node[font=\scriptsize\sffamily, fill=white, inner sep=1.5pt, anchor=west,
      text=FRMGreen] at (axis cs:14.12,7.95) {optimal portfolio, SR $0.551$};
\draw[-{Stealth[length=4pt]}, FRMGreen, thin] (axis cs:14.5,7.92) -- (axis cs:14.02,7.74);
\end{axis}
\end{tikzpicture}}%
{The two destinations, and why they are different places. Cutting the euro on
marginal VaR alone lands on the \emph{minimum-risk} portfolio at 13.85\%
volatility. Taking expected return into account as well lands on the
\emph{optimal} portfolio at 13.98\% --- slightly \emph{riskier}, but with a higher
expected return and the steeper ray from the origin, which is the maximum Sharpe
ratio. The minimum-risk portfolio's expected return of 7.56\% is computed from its
own weights, $0.8521(8\%) + 0.1479(5\%)$, and its Sharpe ratio of 0.546 is lower
than the optimum's 0.551. Minimising risk is not the same as optimising, and the
gap between the two marked points is the difference between risk management and
portfolio management.}

% ------------------------------------------------------------------
\lo{IM-5 e}{Explain the difference between risk management and portfolio management and describe how to use marginal VaR in portfolio management.}

\begin{imdefbox}[title={The difference in one line}]
\term{Risk management} asks only how to make the portfolio less risky, and its
answer is the global minimum where all marginal VaRs are equal.
\term{Portfolio management} asks how to get the best combination of expected
return \emph{and} risk, and its answer is the optimal portfolio where all
excess-return-to-marginal-VaR ratios are equal.
\end{imdefbox}

Marginal VaR is the instrument in both cases. What changes is what it is divided
by: nothing in the first case, and the position's excess expected return in the
second.

\begin{imtrapbox}[title={Consolidated trap summary --- IM-5}]
\begin{itemize}
\item \term{Notation.} $w_i$ is a weight, $W$ or $P$ is portfolio value, $P_i$ or
$x_i$ is a dollar position. The source notes use $w$ for two of these three and
$W$ for a position. Read what the symbol stands for in the stem, every time.
\item \term{Intermediate result, IM-5 a.} Five positions of \$2 million is a \$10
million portfolio. Individual VaR is offered when component VaR is asked for; the
portfolio's total value is often a red herring.
\item \term{Sibling, IM-5 a.} Marginal VaR is a derivative (one more dollar);
incremental VaR is a whole trade; component VaR partitions what you already hold.
Only component VaRs sum to portfolio VaR.
\item \term{Approximation, IM-5 c.} Component VaR $\ne$ incremental VaR. The
first travels down the tangent, the second down the curve. They agree only when
the position is small relative to the portfolio.
\item \term{Formula, IM-5 a.} $\text{CVaR}_i = \text{VaR}_i \times \rho_i$ and
$\text{MVaR}_i = (\text{VaR}_p/P)\beta_i$. Given a correlation, reach for the
first; given a beta, the second.
\item \term{Polarity, IM-5 d.} Cut the position with the \emph{highest} marginal
VaR; add to the one with the \emph{lowest}. For the optimal portfolio it reverses:
add to the \emph{highest} excess-return-to-marginal-VaR ratio.
\item \term{Sibling, IM-5 d.} Minimum-risk is not optimal. The optimal portfolio
here is \emph{riskier} than the minimum-risk one, at 13.98\% against 13.85\%, and
that is correct.
\item \term{Definition, IM-5 b.} Diversification has a floor at
$\sigma\sqrt{\rho}$, not at zero. Only $\rho = 0$ decays to nothing.
\item \term{Units, IM-5 a.} When the position vector is in millions, the resulting
$\sigma_P$ is in millions of dollars, not a percentage. Do not multiply by the
portfolio value a second time.
\end{itemize}
\end{imtrapbox}

% ==================================================================
